\title{Chain-of-Thought Prompting Elicits Reasoning in Large Language Models}

\begin{abstract}
We investigate how producing a \textit{chain of thought}—a sequence of intermediate reasoning steps—substantially enhances the capacity of large language models to carry out complex reasoning. Specifically, we demonstrate that such reasoning skills emerge naturally in sufficiently large language models through a straightforward technique called \textit{chain-of-thought prompting}, where a few examples of chains of thought are included as exemplars in the prompt.

Tests conducted on three large language models reveal that chain-of-thought prompting boosts performance across various arithmetic, commonsense, and symbolic reasoning challenges. The observed improvements can be remarkable. For example, prompting PaLM 540B with only eight chain-of-thought examples achieves state-of-the-art results on the GSM8K dataset of math word problems, outperforming even a finetuned GPT-3 with an additional verifier.
\end{abstract}

\section{Introduction}
The field of NLP has recently undergone a transformation driven by language models \citep[][\textit{inter alia}]{peters-etal-2018-deep,devlin-etal-2019-bert,brown2020language}. Increasing the scale of language models has been shown to bring numerous advantages, such as enhanced performance and better sample efficiency \citep[\textit{inter alia}]{kaplan2020scaling,brown2020language}. Nonetheless, merely enlarging model size has not been enough to achieve high accuracy on difficult tasks like arithmetic, commonsense, and symbolic reasoning \citep{rae2021scaling}.

In this study, we examine how the reasoning capability of large language models can be unlocked using a simple approach inspired by two principles. First, methods for arithmetic reasoning benefit from producing natural language explanations that lead up to the final solution. Previous work has enabled models to produce natural language intermediate steps either by training from scratch \citep{ling-etal-2017-program} or by finetuning pretrained models \citep{cobbe2021training}, as well as neuro-symbolic approaches employing formal languages instead of natural language \citep{roy-roth-2015-solving, chiang-chen-2019-semantically,amini-etal-2019-mathqa, chen2019neural}. Second, large language models introduce the compelling possibility of in-context few-shot learning through \emph{prompting}. Rather than finetuning a new model checkpoint for each task, one can simply prompt the model with a few input-output examples illustrating the task. Notably, this approach has succeeded across a variety of straightforward question-answering tasks \citep{brown2020language}.

Both of the aforementioned concepts, however, have significant drawbacks. For rationale-augmented training and fine-tuning techniques, generating a substantial collection of high-quality rationales is expensive and far more complex than the straightforward input--output pairs typically used in conventional machine learning. The classic few-shot prompting approach employed in \citet{brown2020language} performs poorly on tasks demanding reasoning skills and often shows limited improvement even as the language model size increases \citep{rae2021scaling}. In this study, we merge the advantages of these two strategies in a manner that circumvents their respective weaknesses. Specifically, we investigate the capability of language models to carry out few-shot prompting on reasoning tasks, using prompts composed of triples: $\langle$input, \emph{chain of thought}, output$\rangle$.

A \emph{chain of thought} refers to a sequence of intermediate natural language reasoning steps that culminate in the final output, and we denote this method as \textit{chain-of-thought prompting}. An illustrative example prompt is provided in \cref{fig:pull-figure}.

We conduct empirical assessments on benchmarks involving arithmetic, commonsense, and symbolic reasoning, demonstrating that chain-of-thought prompting surpasses conventional prompting, sometimes by a remarkable margin. \cref{fig:pull-bar-chart} presents one such outcome—on the GSM8K dataset of math word problems \citep{cobbe2021training}, chain-of-thought prompting with \palm{} 540B significantly outperforms standard prompting and establishes new state-of-the-art results. The importance of a prompting-only method lies in its lack of reliance on a large training dataset and the capability of a single model checkpoint to handle numerous tasks without compromising generality. This work highlights how large language models can acquire knowledge from only a few examples using natural language task data (as opposed to learning the underlying input-output patterns automatically through extensive training data).

\section{Chain-of-Thought Prompting} When tackling a complex reasoning problem, such as a multi-step math word problem, people often break the problem down into smaller steps and address each sequentially before arriving at the final solution: \textit{``After Jane gives 2 flowers to her mom she has 10 $\ldots$ then after she gives 3 to her dad she will have 7 $\ldots$ so the answer is 7.''} The objective of this work is to enable language models to produce a similar \textit{chain of thought}—a logical sequence of intermediate reasoning steps that culminate in the final answer. We demonstrate that language models of sufficient scale can generate these chains of thought when examples illustrating chain-of-thought reasoning are included in few-shot prompting.

\cref{fig:pull-figure} presents an instance where a model generates a chain of thought to correctly solve a math word problem it would otherwise fail. The chain of thought here resembles a solution and can be viewed as such, but we choose to refer to it as a chain of thought to emphasize that it imitates a stepwise reasoning process leading to the answer (in contrast to solutions or explanations which generally follow the final answer \citep[][\textit{inter alia}]{narang2020wt5,wiegreffe2021reframing,lampinen2022can}).

Chain-of-thought prompting offers multiple appealing advantages as a method for enhancing reasoning abilities in language models.

\begin{enumerate}[topsep=1pt,itemsep=0ex]
    \item Firstly, chain of thought conceptually enables models to break down complex problems into sequential intermediate steps, allowing for additional computation to be dedicated to problems that involve more stages of reasoning.
    \item Secondly, a chain of thought offers a transparent view into the model’s decision-making process, indicating how it may have derived a specific answer and facilitating debugging when the reasoning path contains errors (although fully understanding the internal computations that lead to an answer remains a challenge).
    \item Thirdly, chain-of-thought reasoning can be applied to tasks such as mathematical word problems, commonsense reasoning, and symbolic manipulation, and it is potentially applicable (at least theoretically) to any problem solvable by humans through language.
    \item Lastly, chain-of-thought reasoning can be easily prompted in sufficiently large pre-trained language models simply by including chain-of-thought examples as part of the few-shot prompt exemplars.
\end{enumerate}

Through empirical studies, we will demonstrate the effectiveness of chain-of-thought prompting on arithmetic reasoning (\cref{sec:arithmetic-reasoning}), commonsense reasoning (\cref{sec:commonsense-reasoning}), and symbolic reasoning (\cref{sec:symbolic-reasoning}).

\section{Arithmetic Reasoning}\label{sec:arithmetic-reasoning}
We start by examining math word problems as depicted in \cref{fig:pull-figure}, which serve to evaluate the arithmetic reasoning capability of language models. Although straightforward for humans, arithmetic reasoning remains a challenging task for language models \citep[][\textit{inter alia}]{hendrycks2021measuring,patel-etal-2021-nlp}. Remarkably, chain-of-thought prompting with the 540B parameter model achieves performance on par with models specifically fine-tuned for these tasks, even setting new state-of-the-art results on the demanding GSM8K benchmark \citep{cobbe2021training}.

\subsection{Experimental Setup}

We investigate the use of chain-of-thought prompting across several language models and multiple evaluation benchmarks.

\textbf{Benchmarks.}
We evaluate on five math word problem datasets:
\textbf{(1)} the \textbf{GSM8K} benchmark of math word problems \citep{cobbe2021training},
\textbf{(2)} the \textbf{SVAMP} dataset consisting of math problems with varied structures \citep{patel-etal-2021-nlp},
\textbf{(3)} the \textbf{ASDiv} dataset covering a diverse range of math word problems \citep{miao-etal-2020-diverse},
\textbf{(4)} the \textbf{AQuA} dataset focusing on algebraic word problems, and
\textbf{(5)} the \textbf{MAWPS} benchmark \citep{koncel-kedziorski-etal-2016-mawps}. Examples from these datasets are provided in Appendix \cref{tab:math-datasets}.

\textbf{Standard prompting.} For the baseline condition, we use conventional few-shot prompting as introduced by \citet{brown2020language}, where a language model is provided with several in-context examples pairing inputs with outputs prior to making predictions on test instances. These exemplars are formatted as question-answer pairs, with the model generating the answer directly, as illustrated in \cref{fig:pull-figure} (left).

\textbf{Chain-of-thought prompting.} Our method involves enhancing each example in few-shot prompting by including a chain of thought that leads to the corresponding answer, as depicted in \cref{fig:pull-figure} (right). Since most datasets only provide an evaluation split, we manually created a set of eight few-shot exemplars with chains of thought for use in prompting—\cref{fig:pull-figure} (right) displays one such exemplar, and the complete set is available in Appendix \cref{tab:appendix-mwp-prompts}. (These exemplars were not subject to prompt engineering; the stability of this approach is examined in \cref{subsec:robustness} and \cref{subsec:faq-prompt-engineering}.) To test whether this style of chain-of-thought prompting can effectively trigger accurate reasoning across diverse math word problems, we utilized this single set of eight chain-of-thought exemplars for all benchmarks except AQuA, which is a multiple-choice task rather than a free-response one. For AQuA, we employed four exemplars and solutions drawn from the training set, as detailed in Appendix \cref{tab:appendix-aqua-prompt}.

\textbf{Language models.} We assess five large language models. The first is \textbf{GPT-3} \citep{brown2020language}, where we utilize the variants text-ada-001, text-babbage-001, text-curie-001, and text-davinci-002, which likely correspond to InstructGPT models with 350M, 1.3B, 6.7B, and 175B parameters respectively \citep{ouyang2022training}. The second is \textbf{\lamda{}} \citep{thoppilan2022lamda}, available in sizes of 422M, 2B, 8B, 68B, and 137B parameters. The third is \textbf{\palm{}}, with model sizes of 8B, 62B, and 540B parameters. The fourth is \textbf{UL2 20B} \citep{tay2022unifying}, and the fifth is \textbf{Codex} \citep[][code-davinci-002 in the OpenAI API]{chen2021evaluating}. For generating outputs from these models, we employ greedy decoding (noting that subsequent research indicates chain-of-thought prompting may be improved by taking the majority vote among multiple sampled generations \citep{wang2022self}). For \lamda{}, we present results averaged over five different random seeds, each corresponding to a unique random ordering of exemplars. Since \lamda{} showed limited variance across these seeds, to conserve computational resources we report results using a single exemplar order for all other models.

\subsection{Results}
The most significant findings regarding chain-of-thought prompting are summarized in \cref{fig:main-math}, with complete experimental outputs for each model collection, model size, and benchmark detailed in \cref{tab:all-lm-math} in the Appendix. There are three main insights. First, \cref{fig:main-math} demonstrates that chain-of-thought prompting emerges as a capability only at larger model scales \citep{wei2022emergent}. Specifically, chain-of-thought prompting does not enhance performance for smaller models and only results in performance improvements when applied to models with around 100B parameters. Qualitative analysis revealed that smaller models generated fluent but illogical reasoning steps, causing worse performance compared to standard prompting.

\input{main_fables/main-math}
Second, chain-of-thought prompting produces greater performance improvements on more complex problems. For example, on GSM8K (which has the lowest baseline performance), performance more than doubled for the largest GPT and \palm{} models. In contrast, for SingleOp, the simplest subset of MAWPS requiring only a single-step solution, performance gains were either negligible or negative (see Appendix \cref{tab:mawps-subsets}).

Third, chain-of-thought prompting with GPT-3 175B and \palm{} 540B shows competitive results relative to previous state-of-the-art approaches, which generally fine-tune task-specific models on labeled training datasets. \cref{fig:main-math} illustrates how \palm{} 540B leverages chain-of-thought prompting to establish new state-of-the-art performance on GSM8K, SVAMP, and MAWPS (noting that standard prompting had already surpassed the previous best on SVAMP). For the other two datasets, AQuA and ASDiv, \palm{} with chain-of-thought prompting achieves results within 2\% of the state of the art (Appendix \cref{tab:all-lm-math}).

To gain a deeper understanding of why chain-of-thought prompting is effective, we conducted a manual review of the reasoning chains generated by \lamda{} 137B for GSM8K. Among 50 randomly selected instances where the model produced the correct final answer, nearly all the generated reasoning chains were logically and mathematically accurate, with only two exceptions that coincidentally led to the correct answer (refer to \cref{subsec:correct-chain-of-thought} and \cref{tab:appendix-gsm8k-correct-examples} for examples of accurate model-generated reasoning chains). Additionally, we examined 50 random cases where the model provided incorrect answers. Our analysis revealed that 46\% of these reasoning chains were almost correct, suffering only from minor errors (such as calculator mistakes, symbol mapping issues, or a single missing reasoning step), while the remaining 54\% contained significant errors in semantic understanding or logical coherence (see \cref{subsec:incorrect-chain-of-thought-analysis}). 

To offer some insight into how scaling enhances chain-of-thought reasoning capabilities, we conducted a similar error analysis on \palm{} 62B and investigated whether scaling up to \palm{} 540B resolved these errors. Our findings indicate that scaling \palm{} to 540B substantially reduces the frequency of one-step omissions and semantic understanding errors present in the 62B model (see \cref{subsec:why-scale-helps}).

\subsection{Ablation Study}\label{subsec:math-ablation}

The demonstrated advantages of chain-of-thought prompting naturally lead to the question of whether similar performance gains can be achieved through alternative forms of prompting. 
\cref{fig:bar_ablation} presents an ablation study exploring three variants of chain-of-thought prompting, which are described below.

\textbf{Equation only.} One hypothesis for why chain-of-thought prompting is beneficial is that it produces the mathematical equation to be solved. To test this, we experiment with a variation where the model is prompted to generate only a mathematical equation prior to providing the final answer. 
As shown in \cref{fig:bar_ablation}, this equation-only prompting approach offers little improvement for GSM8K, suggesting that the semantic complexity of GSM8K questions makes direct translation into equations difficult without the intermediate natural language reasoning steps provided by chain-of-thought. However, for datasets featuring one-step or two-step problems, we observe that equation-only prompting does enhance performance, as the equation can be readily extracted from the question (see Appendix \cref{tab:ablations-arithmetic}).

\input{main_fables/ablation-bar}
\textbf{Variable compute only.} One hypothesis is that chain of thought enables the model to allocate more computation (i.e., generate more intermediate tokens) when dealing with more challenging problems. To separate the impact of variable computation from that of chain-of-thought reasoning, we evaluate a setting where the model is prompted to produce only a sequence of dots ($\ldots$), with the count matching the number of characters in the equation needed to solve the problem. This variant shows performance roughly equivalent to the baseline, indicating that variable computation alone does not account for the effectiveness of chain-of-thought prompting, and that expressing intermediate steps in natural language appears to provide added value.

\textbf{Chain of thought after answer.} Another possible advantage of chain-of-thought prompting might be that it helps the model better access relevant knowledge obtained during pretraining. To examine this, we try a configuration where the chain of thought prompt is presented only after the answer, which helps determine whether the model relies on the generated chain of thought to produce the final answer. This variant achieves performance comparable to the baseline, implying that the sequential reasoning represented by the chain of thought offers benefits beyond merely activating stored knowledge.

\input{main_fables/main-robustness}
\subsection{Robustness of Chain of Thought}\label{subsec:robustness}

Sensitivity to exemplar selection is a critical issue in prompting methods—for example, changing the order of few-shot exemplars can cause GPT-3’s accuracy on SST-2 to vary from near random guessing (54.3\%) to nearly state-of-the-art performance (93.4\%) \citep{zhao2021calibrate}. In this final subsection, we assess the robustness of chains of thought authored by different annotators. Besides the results mentioned earlier, which used chains of thought written by Annotator A, two additional co-authors of this paper (Annotators B and C) independently generated chains of thought for the same few-shot exemplars (detailed in \cref{sec:appendix-alternate-annotators}). Annotator A also produced a more concise version of the chain of thought, following the style of solutions in \citet{cobbe2021training}.\footnote{For example, while the original chain of thought employs several brief sentences (\textit{``There were originally 9 computers. For each of 4 days, 5 more computers were added. So 5 * 4 = 20 computers were added. 9 + 20 is 29.''}), the concise version states \textit{``5 * 4 = 20 new computers were added. So there are 9 + 20 = 29 new computers in the server room now''}.}

\cref{fig:bar-robustness-analysis} presents these findings for \lamda{} 137B on GSM8K and MAWPS (ablation results for other datasets are provided in Appendix \cref{tab:ablations-arithmetic} / \cref{tab:ablations-commonsense-symbolic}). While there is variability across different chain of thought annotations, as anticipated when employing exemplar-based prompting \citep{le-scao-rush-2021-many,reynolds2021prompt,zhao2021calibrate}, all sets of chain of thought prompts significantly outperform the standard baseline. This indicates that effective use of chain of thought does not rely on any specific linguistic style.

To verify that effective chain-of-thought prompting generalizes to other exemplar sets, we also conducted experiments with three sets of eight exemplars randomly selected from the GSM8K training set, which is an independent source (examples in this dataset already contain reasoning steps resembling a chain of thought).\footnote{We selected examples $\leq 60$ tokens to fit within our input context window and restricted the examples to $\leq 2$ steps to ensure a fair comparison with the eight manually composed exemplars.} \cref{fig:bar-robustness-analysis} demonstrates that these prompts performed similarly to our manually crafted exemplars, also greatly surpassing standard prompting.

Beyond robustness to annotator variability, independently authored chains of thought, different exemplars, and various language models, we also observe that chain-of-thought prompting for arithmetic reasoning is resilient to different exemplar orders and diverse numbers of exemplars (refer to \cref{subsec:faq-prompt-engineering}).

\section{Commonsense Reasoning}\label{sec:commonsense-reasoning}

Although chain of thought is especially well-suited for math word problems, its language-based nature makes it applicable to a wide range of commonsense reasoning tasks, which entail reasoning about physical and human interactions based on general background knowledge. Commonsense reasoning is crucial for interacting with the world and remains beyond the capabilities of current natural language understanding systems \citep{talmor2022commonsenseqa}.

\textbf{Benchmarks.}  
We utilize five datasets that encompass a wide variety of commonsense reasoning challenges. The well-known \textbf{CSQA} \citep{talmor-etal-2019-commonsenseqa} presents commonsense questions about the world, often involving complex semantics that typically necessitate background knowledge.  
\textbf{StrategyQA} \citep{geva-etal-2021-aristotle} requires models to deduce a multi-step strategy in order to answer the questions. We select two specialized evaluation datasets from the BIG-bench project \citep{bigbench}: \textbf{Date} Understanding, which requires inferring dates from provided contexts, and \textbf{Sports} Understanding, which involves judging the plausibility or implausibility of sentences related to sports. Lastly, the \textbf{SayCan} dataset \citep{ahn2022can} tasks models with translating natural language instructions into sequences of robot actions chosen from a discrete set.  
\cref{fig:dataset-examples} presents sample instances with chain of thought annotations for all datasets.

\textbf{Prompts.}  
We adopt the same experimental protocol as in the previous section. For CSQA and StrategyQA, we randomly picked examples from the training set and manually crafted chains of thought to serve as few-shot exemplars. Since the two BIG-bench tasks lack training sets, we selected the first ten examples from the evaluation set as few-shot exemplars and evaluated on the remaining examples. For SayCan, we used six examples from the training set utilized by \citet{ahn2022can} and also manually constructed chains of thought.

\textbf{Results.} 
\cref{fig:commonsense-results} presents these findings specifically for \palm{} (complete results for \lamda{}, GPT-3, and various model sizes are provided in \cref{tab:all-lm-commonsense}). 
Across all tasks, increasing the model size enhanced the effectiveness of standard prompting; adding chain-of-thought prompting yielded additional improvements, with the largest gains observed for \palm{} 540B. When using chain-of-thought prompting, \palm{} 540B demonstrated strong performance relative to baseline methods, surpassing the previous best results on StrategyQA (75.6\% compared to 69.4\%) and exceeding the accuracy of an unaided sports expert in sports understanding (95.4\% versus 84\%). 
These findings indicate that chain-of-thought prompting can also boost performance on tasks that involve diverse commonsense reasoning skills (although improvements were minimal on CSQA).

\input{main_fables/main-commonsense}

\input{main_fables/main-symbolic}
\section{Symbolic Reasoning}\label{sec:symbolic-reasoning} 
Our final set of experiments focuses on symbolic reasoning, which is straightforward for humans but can be difficult for language models. We demonstrate that chain-of-thought prompting not only allows language models to tackle symbolic reasoning tasks that are challenging under standard prompting but also enables generalization to longer input sequences at inference time than those seen in the few-shot examples.

\paragraph{Tasks.} 
We consider the following two simplified tasks.

\begin{itemize}[leftmargin=*,topsep=0pt]
    \itemsep0em \item \textbf{Last letter concatenation.} This task requires the model to join the last letters of words in a name (for example, \textit{``Amy Brown''} $\rightarrow$ \textit{``yn''}). 
    It is a more difficult variant of first letter concatenation, a task language models can already solve without chain of thought.\footnote{We tested 10 common names using GPT-3 \texttt{davinci} and it answered all but one correctly.} Full names are generated by randomly combining names drawn from the top one-thousand first and last names based on name census data ({\small{\url{https://namecensus.com/}}}).
    \item \textbf{Coin flip.} This task asks the model to determine whether a coin remains heads up after a series of people either flip or do not flip it (e.g., \textit{``A coin is heads up. Phoebe flips the coin. Osvaldo does not flip the coin. Is the coin still heads up?''} $\rightarrow$ \textit{``no''}).
\end{itemize}

Given that the design of these symbolic reasoning tasks is clearly specified, for each task we evaluate on an \textit{in-domain} test set where the examples have the same number of steps as the training or few-shot exemplars, as well as on an \textit{out-of-domain} (OOD) test set where the test examples contain more steps than those in the exemplars. For the last letter concatenation task, the model is only exposed to exemplars consisting of two-word names, and is then tested on last letter concatenation for names with 3 and 4 words.\footnote{For names longer than two words, we concatenate multiple first and last names together.} A similar approach is taken for the number of possible flips in the coin flip task. Our experimental framework employs the same methods and models described in the previous two sections. We again manually create chains of thought for the few-shot exemplars corresponding to each task, as presented in \cref{fig:dataset-examples}.

\paragraph{Results.} 
The outcomes of these in-domain and OOD evaluations are displayed in \cref{fig:symbolic-main} for \palm{}, with \lamda{} results provided in Appendix \cref{tab:all-lm-symbolic}. Utilizing \palm{} 540B, chain-of-thought prompting achieves nearly 100\% accuracy (note that standard prompting already solves the coin flip task with \palm{} 540B, though not with \lamda{} 137B). It is important to emphasize that these in-domain evaluations are essentially “toy tasks” because the exact solution procedures are supplied by the chains of thought in the few-shot exemplars; the model merely needs to replicate the same sequence of steps with new symbols during testing. Nevertheless, smaller models continue to fail— the capability to execute abstract manipulations on unseen symbols for these three tasks only emerges at model scales of around 100 billion parameters.

Regarding the OOD evaluations, standard prompting does not succeed for either task. However, with chain-of-thought prompting, language models exhibit upward scaling trends (although their performance remains lower than in the in-domain scenario). Therefore, chain-of-thought prompting supports length generalization beyond previously encountered chains of thought for language models of adequate scale.

\section{Discussion}\label{sec:discussion}

We have investigated chain-of-thought prompting as a straightforward approach to induce multi-step reasoning in large language models. Initially, we observed that chain-of-thought prompting significantly enhances performance on arithmetic reasoning, providing improvements that surpass those seen in ablation studies and remain consistent across different annotators, exemplars, and language models (\cref{sec:arithmetic-reasoning}).
Subsequently, experiments on commonsense reasoning highlighted the general applicability of chain-of-thought reasoning due to its linguistic nature (\cref{sec:commonsense-reasoning}).
Finally, we demonstrated that for symbolic reasoning, chain-of-thought prompting enables OOD generalization to longer sequence lengths (\cref{sec:symbolic-reasoning}).
In all these experiments, chain-of-thought reasoning was elicited merely by prompting a pre-existing language model, with no finetuning applied during this work.

The appearance of chain-of-thought reasoning as a function of model scale has been a consistent pattern \citep{wei2022emergent}. For numerous reasoning tasks where standard prompting exhibits a flat scaling curve, chain-of-thought prompting results in sharply rising scaling trends. It seems that chain-of-thought prompting broadens the spectrum of tasks that large language models can effectively perform—in other words, our findings emphasize that standard prompting represents only a lower bound on the capabilities of large language models. This insight likely provokes more questions than it resolves—for example, to what extent can reasoning abilities further improve as model scale increases? What other prompting strategies might extend the range of tasks language models are capable of solving?

Regarding limitations, we first note that although chain-of-thought mimics the reasoning process of humans, it does not establish whether the neural network is genuinely “reasoning,” which remains an open question. Second, while the manual effort to augment exemplars with chains of thought is minimal in a few-shot setting, the annotation cost could be substantial for finetuning (though this might be alleviated via synthetic data generation or zero-shot generalization).
\orange{Third, there is no assurance that the reasoning paths are correct, which can yield both accurate and inaccurate answers; enhancing the factual accuracy of language model outputs is a promising area for future research \citep[][\textit{inter alia}]{rashkin2021measuring,ye2022unreliability,wiegreffe2021reframing}.} Lastly, the fact that chain-of-thought reasoning emerges only at large model scales makes deployment expensive in practical applications; future work could investigate methods to elicit reasoning in smaller models.

\section{Related Work} This study draws inspiration from various research domains, which we elaborate on in a more comprehensive related work section (\cref{sec:extended-related-work}). Here, we focus on two key directions and the corresponding studies that are arguably the most pertinent.

The first significant direction involves utilizing intermediate steps to address reasoning challenges. \cite{ling-etal-2017-program} introduced the concept of employing natural language explanations to solve math word problems through a sequence of intermediate steps. Their approach stands in notable contrast to research relying on formal languages for reasoning \citep{roy-2016-reasoning, chiang-chen-2019-semantically, amini-etal-2019-mathqa, chen2019neural}. Building on \citet{ling-etal-2017-program}, \citet{cobbe2021training} developed a larger dataset and used it to fine-tune a pretrained language model instead of training a model from scratch. Within the field of program synthesis, \cite{nye2021show} utilize language models to predict the final outputs of Python programs by first predicting intermediate computational results line-by-line, demonstrating that their stepwise prediction approach outperforms directly forecasting the final outputs.

This paper is also closely connected to the extensive recent research on prompting. Since the introduction of few-shot prompting by \citet{brown2020language}, several general methods have enhanced the prompting capabilities of models, such as learning prompts automatically \citep{lester-etal-2021-power} or providing models with task-specific instructions \citep{wei2021finetuned,sanh2021multitask,ouyang2022training}. While these methods focus on improving or supplementing the input portion of the prompt (e.g., adding instructions before inputs), our work takes a complementary approach by augmenting language model outputs with a chain of thought.

\section{Conclusions} We have investigated chain-of-thought prompting as a straightforward and widely applicable technique for improving reasoning in language models. Through experiments involving arithmetic, symbolic, and commonsense reasoning, we observe that chain-of-thought reasoning emerges as a function of model scale, enabling sufficiently large language models to tackle reasoning tasks that otherwise exhibit flat scaling behavior. Expanding the range of reasoning capabilities in language models is expected to encourage further research into language-based reasoning methodologies.

\end{document}